\chapter{Variational Principles }
\label{vp}

\begin{motto} \textbf{Motto.} A variational principle is not a single statement
    but a \emph{family} of statements, distinguished by what is held fixed and
    what is left free at the endpoints. For example, Hamilton fixes time and
    lets energy float; Maupertuis fixes energy and lets time float; Gauss and
    Hertz fix nothing beyond the current instant. The equations of motion are
    common ground; the variational scaffolding around them is not. 
\end{motto}

Classical mechanics is usually presented as though it had one variational
principle: ``the'' principle of least action, $\delta\!\int L\,dt = 0$. In fact
the historical record contains a dozen or so logically distinct variational
statements about mechanical systems. They agree on the trajectories they select
in the generic case, which is why they are folded together in most textbook
treatments; they disagree, instructively, at exactly the points where the
generic case fails --- nonholonomic constraints, free versus fixed transit
time, the logical status of momentum, and the passage to quantum mechanics.
This chapter lays the principles out side by side, proves the equivalences that
hold, and is explicit about the ones that do not.

\begin{definition} Let $Q$ be a configuration space and $q:[t_1,t_2]\to Q$ a
    curve. A \emph{functional} is a map $F$ from curves to $\mathbb{R}$. Given
    a one-parameter family of curves $q_\varepsilon(t)=q(t)+\varepsilon\,\delta
    q(t) + O(\varepsilon^2)$ with $\delta q$ vanishing wherever the endpoints
    are held fixed, the \emph{first variation} is 
    \[ 
    \delta F :=\left.\frac{d}{d\varepsilon}\right|_{\varepsilon=0} F[q_\varepsilon].
    \] 
    A curve is a \emph{stationary point} (extremal) of $F$ if $\delta F=0$ for
    every admissible $\delta q$. 
\end{definition}

For $F[q]=\int_{t_1}^{t_2} L(q,\dot q,t)\,dt$ with fixed endpoints 
$\delta q(t_1)=\delta q(t_2)=0$, integration by parts gives
\[
\delta S = \int_{t_1}^{t_2}\left(\frac{\partial L}{\partial q^i}-
\frac{d}{dt}\frac{\partial L}{\partial \dot q^i}\right)\delta q^i\,dt
+ \left[\frac{\partial L}{\partial \dot q^i}\,\delta q^i\right]_{t_1}^{t_2}.
\]

The boundary (transversality) term is the whole story of this chapter.
Hamilton's principle kills it by fixing $q$ at both ends. Maupertuis' principle
keeps the ends of the \emph{curve} fixed but frees the \emph{time} at which
they are reached, so a different boundary term survives and has to be handled
separately. Gauss's and Hertz's principles dispense with the interval
altogether and vary only at a single instant.

\section{Hamilton's principle}
\label{sec:hamilton}

\begin{definition}
\textbf{Hamilton's principle.} Among curves $q:[t_1,t_2]\to Q$ with 
    \emph{fixed endpoints in time and configuration}, 
    $q(t_1)=q_1$, $q(t_2)=q_2$, $t_1,t_2$ fixed, 
    the physical trajectory renders the action
\[
S[q(\cdot)] = \int_{t_1}^{t_2} L(q,\dot q,t)\,dt
\]
stationary.
\end{definition}

Setting $\delta S=0$ with the boundary term of the previous section 
vanishing gives the Euler--Lagrange equations
\[
\frac{d}{dt}\frac{\partial L}{\partial \dot q^i} 
- \frac{\partial L}{\partial q^i} = 0 .
\]

Hamilton's principle, applied without further thought to a system with
\emph{nonholonomic} velocity constraints $a^{(k)}_i(q)\dot q^i=0$ by adjoining
the constraints to $L$ with Lagrange multipliers before varying, produces the
\emph{vakonomic} equations of motion. These generically disagree with the
correct Lagrange--d'Alembert equations for a nonholonomically constrained
system (the textbook example is a coin or sphere rolling without slipping). The
discrepancy is not a paradox to be explained away: it is a sign that Hamilton's
principle, in its integral, fixed-endpoint form, is simply not the right
variational statement for a constraint that is imposed at the level of
velocities rather than configurations. Section~\ref{sec:gauss} shows which
principle \emph{is} the right one.

\section{Maupertuis' principle}
\label{sec:maupertuis}

\begin{definition}
\label{def:maupertuis-old}
\textbf{Maupertuis' principle (principle of least action, abbreviated form).}
    Suppose $H$ has no explicit time dependence, so $E=H$ is conserved. Among
    curves in configuration space joining $q_1$ to $q_2$ that are compatible
    with energy $E$ --- with the \emph{time of transit left free} --- the
    physical trajectory renders the abbreviated action
\[
W[q(\cdot)] = \int p_i\,dq^i
\]
stationary, where along each trial curve $p_i$ is fixed pointwise by the energy
    constraint $H(q,p)=E$.
\end{definition}

\begin{remark}[Jacobi's complaint]
\label{rem:jacobi-complaint}
Definition~\ref{def:maupertuis-old} constrains the energy \emph{pointwise} on every trial curve, $H(q,p)=E$, not merely on the true trajectory. Jacobi  objected, that the principle is presented in essentially every textbook so as to be nearly unintelligible, and the pointwise constraint is the reason why: it presupposes on trial curves the very energy conservation that Hamilton's principle only \emph{derives} for the true one, and in low dimension it leaves almost nothing left to vary (fixing $E=v^2/2$ for a free particle on a line, for instance, fixes the speed outright). 
\end{remark}

\begin{theorem}[equivalence, isoenergetic case]
Among curves of fixed energy $E$ joining $q_1$ to $q_2$, with transit time
    free, $\delta W=0$ if and only if the curve, traversed according to $H=E$,
    satisfies Hamilton's equations.
\end{theorem}

$W$ and $S$ therefore extremize over \emph{different} classes of curves and are
not interchangeable term by term: $S$'s stationarity theorem fixes the
parametrization (time) and lets energy be whatever it turns out to be; $W$'s
fixes the energy and is blind to parametrization, recovering only the
trajectory as a set of points in $Q$ (the orbit shape), with the
time-parametrization reconstructed afterward from $H=E$. This is why
Maupertuis' principle is the natural one for questions about orbit \emph{shape}
(is the Kepler orbit closed? what is its precession?) while Hamilton's is the
natural one for questions about time evolution.

\section{Reciprocal principles and the generalized Maupertuis principle}
\label{sec:reciprocal}

Gray, Karl and Novikov repaired Maupertuis' principle by weakening the
pointwise energy constraint of Definition~\ref{def:maupertuis-old} to a
constraint on the trial curve's \emph{time-averaged} energy alone --- a
constraint every trial curve can satisfy, and one that no longer presupposes
what Hamilton's principle is supposed to derive.

For a trial trajectory $q(\cdot)$ of duration $T=t_2-t_1$ joining fixed
    $q_1,q_2$, the \emph{mean energy} is
\[
\bar E := \frac1T\int_{t_1}^{t_2} H(q,p,t)\,dt .
\]
On the true trajectory of a conservative system $\bar E=E$, but on an arbitrary
    trial curve $\bar E$ is perfectly well defined even when $H$ is not
    pointwise constant along it.

\begin{definition}
\textbf{The Generalized Maupertuis Principle (GMP).} Among trial curves with
    fixed endpoints $q_1,q_2$ and fixed \emph{mean} energy $\bar E$, transit
    time $T$ left free, the true trajectory renders $W=\int p_i\,dq^i$
    stationary:
\[
(\delta W)_{\bar E} = 0 .
\]
\end{definition}

Varying this identity at fixed endpoints, with both the trial curve and the
duration $T$ free, splits into two independent statements: fixing $T$ recovers
Hamilton's principle $(\delta S)_T=0$; fixing $\bar E$ instead gives the GMP
$(\delta W)_{\bar E}=0$. The two are therefore equivalent to each other, and
--- unlike the textbook Maupertuis principle of
Definition~\ref{def:maupertuis-old} --- the GMP derives energy conservation for
the true trajectory rather than assuming it, resolving Jacobi's complaint.

\begin{definition}
Reversing which quantity is held fixed and which is left free in each principle
    gives its \emph{reciprocal}: the \textbf{Reciprocal Maupertuis Principle
    (RMP)},
\[
(\delta \bar E)_W = 0,
\]
and the \textbf{Reciprocal Hamilton Principle (RHP)},
\[
(\delta T)_S = 0 .
\]
The RMP is a principle of \emph{least (stationary) mean energy} at fixed action
    $W$; the RHP is a principle of \emph{least time} at fixed action $S$.
\end{definition}

\begin{theorem}[the reciprocity--Legendre square]
The four statements $(\delta \bar E)_W=0$, $(\delta W)_{\bar E}=0$, $(\delta
    S)_T=0$, and $(\delta T)_S=0$ are pairwise equivalent (each characterizes
    the same true trajectories): horizontal neighbors in
    Figure~\ref{fig:square} are related by \emph{reciprocity} (swapping which
    member of a conjugate pair is fixed and which is varied), and vertical
    neighbors by the \emph{Legendre} relation $\bar ET=W-S$.
\end{theorem}

\begin{figure}[h]
\centering
\begin{tikzpicture}[
  box/.style={draw, rounded corners, align=center, minimum height=8mm, minimum width=30mm, inner sep=5pt, font=\small},
  arr/.style={thick, {Latex[length=2mm]}-{Latex[length=2mm]}},
  lbl/.style={font=\scriptsize, fill=white, inner sep=1pt}
]
  \node[box, fill=blue!8] (rmp) at (0,2.1) {RMP\\ $(\delta \bar E)_W=0$};
  \node[box, fill=blue!8] (gmp) at (5,2.1) {GMP\\ $(\delta W)_{\bar E}=0$};
  \node[box, fill=blue!8] (rhp) at (0,0) {RHP\\ $(\delta T)_S=0$};
  \node[box, fill=blue!8] (hp) at (5,0) {HP\\ $(\delta S)_T=0$};
  \draw[arr] (rmp) -- node[lbl,above]{reciprocity} (gmp);
  \draw[arr] (rhp) -- node[lbl,below]{reciprocity} (hp);
  \draw[arr] (rmp) -- node[lbl,left]{Legendre} (rhp);
  \draw[arr] (gmp) -- node[lbl,right]{Legendre} (hp);
\end{tikzpicture}
\caption{The reciprocity--Legendre square of Gray, Karl and Novikov. Horizontal arrows swap what is held fixed and what is varied; vertical arrows exchange $(W,\bar E)$ for $(S,T)$ via $\bar ET=W-S$.}
\label{fig:square}
\end{figure}

\section{Jacobi's principle: geometrization}
\label{sec:jacobi}

For a natural system, $T=\tfrac12 g_{ij}(q)\dot q^i\dot q^j$ defines a
Riemannian metric $g_{ij}$ on $Q$ (the mass metric), and $p_i\,dq^i =
g_{ij}\dot q^j\,dq^i = 2T\,dt$ along any trajectory. Writing
$ds^2=g_{ij}\,dq^idq^j$ for the mass-metric line element, $dt =
ds/\sqrt{2T}=ds/\sqrt{2(E-V)}$, so
\[
W=\int 2(E-V)\,dt = \int \sqrt{2(E-V)}\;ds .
\]

\begin{definition}
\label{def:jacobi-metric}
The \emph{Jacobi metric} is the conformal rescaling
\[
ds_J^2 := 2\big(E-V(q)\big)\,g_{ij}(q)\,dq^i\,dq^j,
\]
defined on the region $\{E>V\}$ where motion is classically allowed.
\end{definition}
Maupertuis' principle is exactly the statement that trajectories of energy $E$
are \emph{geodesics} of the Jacobi metric: $\delta\!\int ds_J = 0$.

This is Jacobi's reformulation: dynamics at fixed energy is pure Riemannian
geometry, with the potential absorbed into the metric. Two consequences are
worth flagging because they recur elsewhere in this book: (i) the curvature of
$g^J_{ij}=2(E-V)g_{ij}$ controls the stability of nearby trajectories exactly
as sectional curvature controls geodesic deviation, giving a direct route from
a potential $V(q)$ to exponential trajectory divergence (Krylov's original
argument for mixing in the hard-sphere gas, later formalized for Anosov
geodesic flows); (ii) since $Q$ and its metric are usually finite-dimensional
and time has been eliminated entirely, the Jacobi picture is the cleanest
illustration that ``geometrize first, quantize/perturb second'' is a genuine
choice of route, not a cosmetic rewriting.

\section{A genuine non-equivalence}
\label{sec:nonequivalence}

Every equivalence proved so far ---  the isoenergetic theorem, the
reciprocity--Legendre square of Section~\ref{sec:reciprocal} --- is
\emph{local}: each relates a single given trajectory to itself, under different
fixed/free bookkeeping. None of them says that the \emph{set} of trajectories
obtained by prescribing $T$ in advance (Hamilton) coincides with the set
obtained by prescribing $E$ in advance (Maupertuis/GMP). Once the potential has
more than one local minimum, these two sets --- and even their cardinalities
--- generically differ.

Hamilton's principle imposes no restriction on trial curves at all: $q(t)$ may
pass through any point of $Q$, at any speed, en route from $q_1$ to $q_2$,
since $L=T(\dot q)-V(q)$ is perfectly well defined however fast or slow a trial
curve moves at a given $q$. The textbook Maupertuis principle
(Definition~\ref{def:maupertuis-old}) is different in kind: $p$ along the trial
curve is \emph{defined} by $H(q,p)=E$, so $T=E-V(q)$ must be $\ge0$ pointwise
--- a trial curve can only exist within the classically allowed (Hill's) region
$Q_E:=\{q\in Q: V(q)\le E\}$, exactly the domain restriction already flagged in
the Jacobi metric of Definition~\ref{def:jacobi-metric}. If $q_1$ and $q_2$ lie
in different path-components of $Q_E$ --- separated by a barrier taller than
$E$ --- the textbook Maupertuis competition at that $E$ has no trial curves
joining them at all, not merely no minimizer. The GMP of
Section~\ref{sec:reciprocal} repairs this too: since only the \emph{mean}
energy is fixed, a GMP trial curve is exactly as free as a Hamilton one,
dipping in and out of $\{V>E\}$ as it pleases, provided the time-average works
out to $\bar E$. But this freedom is only in the trial family: the actual
extremal, once found, is a genuine solution of an autonomous Hamiltonian system
and so has \emph{constant} energy equal to its own mean. If that constant falls
below a barrier separating $q_1$ from $q_2$, no connecting extremal exists in
the GMP formulation any more than in Hamilton's --- not because of how the
competition was set up, but because ordinary energy conservation forbids it.

Even granting accessibility, the \emph{number} of genuine trajectories
connecting two points at one fixed $E$ need not match the number connecting
them in one fixed $T$. Let $\Omega_E:=\{q:V(q)\le E\}$ be a bounded potential
well with $E$ a regular value of $V$. Via the Jacobi metric, a trajectory that
oscillates back and forth within $\Omega_E$ --- a \emph{brake orbit}, so called
because the particle momentarily comes to rest, $p=0$, at each end of its
swing, on $\partial\Omega_E$ --- corresponds exactly to a geodesic chord of
$ds_J$ meeting $\partial\Omega_E$ orthogonally. Seifert (1948) proved that at
least one such orbit always exists, and conjectured at least $N=\dim Q$
geometrically distinct ones whenever $\Omega_E$ is homeomorphic to an $N$-disk:
genuine multiplicity of fixed-energy solutions, expected already in a single,
topologically trivial well, before a second minimum is ever introduced. The
conjecture is now a theorem under a concavity hypothesis on $\partial\Omega_E$
(Giambò, Giannoni and Piccione), and is known to fail without some such
hypothesis: the same authors exhibit real-analytic wells, homeomorphic to a
disk, carrying only \emph{one} brake orbit. None of this multiplicity is
visible from the Hamilton side by fixing a single $T$ and solving once: each
brake orbit has its own period, generically different from the others, so a
fixed-$T$ search catches at most whichever one happens to match.

None of this contradicts Section~\ref{sec:reciprocal}: trajectory by
trajectory, Hamilton's and Maupertuis'/GMP's principles remain exactly
equivalent. What fails is the stronger, unstated assumption that fixing $T$ and
fixing $E$ pose ``the same problem'' with the same solution count. Once
$\Omega_E$ has nontrivial topology --- exactly what a potential with several
minima produces, once $E$ rises above the barriers separating them --- fixing
$E$ and fixing $T$ become genuinely different slicings of one underlying family
of solutions: governed, on the Maupertuis side, by the topology of $\Omega_E$
(Seifert, Ljusternik--Schnirelmann category); and, on the Hamilton side, by a
related but distinct Morse theory of the fixed-time path space. ``Maupertuis'
principle is equivalent to Hamilton's principle'' is therefore true of
individual trajectories, and false, in general, of the two boundary-value
\emph{problems}.

\section{Gauss's principle of least constraint}
\label{sec:gauss}

Hamilton's and Maupertuis' principles are \emph{integral} statements: they
compare an entire trial curve against neighboring curves over a finite stretch
of time or a finite arc. Gauss's principle is of a different logical kind: it
is \emph{differential}, comparing only candidate accelerations at the current
instant.

\begin{definition}
\textbf{Gauss's principle of least constraint.} For particles of mass $m_a$ at
    given positions and velocities, subject to (possibly nonholonomic)
    constraints, and acted on by applied forces $\vec F_a$, the true
    accelerations $\ddot{\vec r}_a$ minimize
\[
Z(\ddot{\vec r}) := \sum_a m_a\left|\ddot{\vec r}_a 
    - \frac{\vec F_a}{m_a}\right|^2
\]
over all accelerations compatible with the constraints differentiated to the
    acceleration level, positions and velocities held fixed.
\end{definition}

Gauss's principle is equivalent to d'Alembert's principle (and hence to
Newton's second law together with the constraint forces doing no virtual work).
Minimizing the quadratic form $Z$ over the affine space of admissible
$\ddot{\vec r}_a$ sets its gradient orthogonal (in the mass inner product) to
the tangent space of admissible \emph{variations} $\delta\ddot{\vec r}_a$ of
that space, i.e.
\[
\sum_a \big(m_a\ddot{\vec r}_a - \vec F_a\big)\cdot \delta\ddot{\vec r}_a = 0
\]
for every admissible $\delta \ddot{\vec r}_a$. For the constraints in question,
admissible acceleration-variations and admissible virtual displacements $\delta
\vec r_a$ span the same subspace, so this is exactly d'Alembert's principle,
$\sum_a(m_a\ddot{\vec r}_a-\vec F_a)\cdot\delta\vec r_a=0$.

Because Gauss's principle is imposed directly on the constraint manifold at the
acceleration level, it reproduces the Lagrange--d'Alembert equations for a
nonholonomic system \emph{by construction}, with no vakonomic detour. The
mismatch flagged in Section~\ref{sec:hamilton} is therefore not a defect in
mechanics but a mismatch of variational \emph{principles} to the constraint
type: integral, fixed-endpoint principles (Hamilton, Maupertuis) are suited to
holonomic constraints; the differential, instantaneous principle (Gauss)
handles nonholonomic ones correctly because it never has to compare curves that
end differently far along a constraint surface that the constraint itself does
not let you parametrize by position alone.

Gauss's principle is usually stated in physical acceleration variables, but it
restates without loss in terms of \emph{quasi-accelerations} $\dot\omega^r$
adapted to the constraints: minimizing $Z$ over admissible quasi-accelerations
gives the \emph{Gibbs--Appell equations} $\partial\mathcal
S/\partial\dot\omega^r = Q_r$, where $\mathcal S:=\tfrac12\sum_a m_a|\ddot{\vec
r}_a|^2$ --- the same quantity Gauss minimizes, re-expressed --- is the
\emph{Gibbs function} and $Q_r$ are the generalized forces conjugate to the
quasi-velocities. Gibbs (1879) and, independently, Appell (1900) introduced
this form because it handles constraints of arbitrarily general type, even
nonlinear in the velocities, without ever introducing a Lagrange multiplier ---
a second, independent route to the resolution Remark~\ref{rem:fragility}
promised.

\section{Hertz's principle of least curvature}
\label{sec:hertz}

\begin{definition}
\textbf{Hertz's principle} A system subject only to holonomic (``rigid'')
    constraints and no applied forces moves, at constant speed, along the path
    of \emph{least curvature} through configuration space equipped with the
    mass metric $g_{ij}$.
\end{definition}

Hertz's principle is the $V\equiv 0$ specialization of Jacobi's principle.
Set $V=0$ in the Jacobi metric: $ds_J^2 = 2E\,g_{ij}\,dq^idq^j$, a
\emph{constant} conformal rescaling of $g_{ij}$. Geodesics are unchanged by a
constant conformal factor (only their affine speed changes, to $\sqrt{2E}$), so
the $ds_J$-geodesics coincide with the $g_{ij}$-geodesics, traversed at
constant speed. ``Least curvature'' is Hertz's own (equivalent) way of
characterizing a geodesic: the path whose curvature vector, computed in
$g_{ij}$, vanishes.

This is precisely the ``forceless mechanics'' discussed elsewhere in this book:
what looks like force-free inertial motion in a curved configuration space
(e.g.\ a bead constrained to a curved wire, or the reduced dynamics of a system
with cyclic coordinates eliminated) is literally geodesic motion once the
constraints have been folded into the metric $g_{ij}$ --- Hertz's proposal to
eliminate force from the foundations of mechanics amounts to insisting that
\emph{all} of dynamics should look like this, with potentials themselves
eventually explained away via hidden constrained (``cyclic'') degrees of
freedom.

\section{Herglotz's principle: a genuine variational principle for dissipation}
\label{sec:herglotz}

Every principle so far --- Hamilton's, Maupertuis'/GMP's, Jacobi's, Gauss's,
Hertz's --- governs conservative, or at worst constrained-but-reversible,
systems: none of them, as stated, can produce a genuinely dissipative equation
of motion, one along which mechanical energy is not conserved. The usual patch,
Rayleigh's dissipation function added by hand to the Euler--Lagrange equations,
is exactly that: a patch, not the output of extremizing any single functional.
Herglotz's principle removes the patch.

\begin{definition}
\textbf{Herglotz's variational principle} (1930). Let $L(q,\dot q,z,t)$ depend
    on an auxiliary variable $z$ as well as on $q,\dot q,t$, and let $z(\cdot)$
    be determined along any trial curve $q(\cdot)$ by
\[
\dot z(t) = L\big(q(t),\dot q(t), z(t), t\big), \qquad z(t_1)=z_1 \text{ fixed}.
\]
The true trajectory renders the \emph{endpoint value} $z(t_2)$ stationary among
    curves with fixed $q(t_1)=q_1$, $q(t_2)=q_2$: $\delta z(t_2) = 0$.
\end{definition}

When $L$ does not depend on $z$, $z(t_2)-z_1=\int_{t_1}^{t_2}L\,dt$ and
Herglotz's principle reduces to Hamilton's. The dependence on $z$ is the whole
extension.

Herglotz's principle is equivalent to the generalized Euler--Lagrange equations
\[
\frac{d}{dt}\frac{\partial L}{\partial\dot q^i} 
- \frac{\partial L}{\partial q^i} 
- \frac{\partial L}{\partial z}\frac{\partial L}{\partial\dot q^i} = 0 .
\]
Indeed, linearize the defining ODE along a variation $\delta q^i$ 
(with $\delta q^i(t_1)=\delta q^i(t_2)=0$, hence 
$\delta z(t_1)=0$): $\delta\dot z =
\frac{\partial L}{\partial q^i}\delta q^i+\frac{\partial L}{\partial \dot
q^i}\delta\dot q^i+\frac{\partial L}{\partial z}\delta z$, a linear ODE for
$\delta z$ along the trajectory. With the integrating factor
$\mu(t):=\exp\!\big(-\int_{t_1}^t \frac{\partial L}{\partial z}\,dt'\big)$, so
$\dot\mu=-\mu\,\partial L/\partial z$, this becomes $\frac{d}{dt}(\mu\,\delta
z) = \mu\left(\frac{\partial L}{\partial q^i}\delta q^i+\frac{\partial
L}{\partial\dot q^i}\delta\dot q^i\right)$. Integrating from $t_1$ to $t_2$ and
integrating the $\delta\dot q^i$ term by parts (the boundary contribution
vanishes since $\delta q^i(t_1)=\delta q^i(t_2)=0$, and
$\frac{d}{dt}\big(\mu\,\partial L/\partial\dot
q^i\big)=\mu\big(\frac{d}{dt}\frac{\partial L}{\partial\dot q^i}-\frac{\partial
L}{\partial z}\frac{\partial L}{\partial\dot q^i}\big)$ since
$\dot\mu/\mu=-\partial L/\partial z$) gives
\[
\mu(t_2)\,\delta z(t_2) = \int_{t_1}^{t_2}\mu\left[\frac{\partial L}{\partial q^i}-
\frac{d}{dt}\frac{\partial L}{\partial\dot q^i}+\frac{\partial L}{\partial z}
\frac{\partial L}{\partial \dot q^i}\right]\delta q^i\,dt .
\]
Since $\mu(t_2)\neq0$, $\delta z(t_2)=0$ for every admissible $\delta q^i$ 
forces the bracket to vanish.

\begin{example}
Take $L=\tfrac12 m\dot q^2 - V(q) - \gamma z$ for constant $\gamma$. Then
    $\partial L/\partial z=-\gamma$, and the generalized Euler--Lagrange
    equation reads
\[
m\ddot q + V'(q) + \gamma m\dot q = 0,
\]
the ordinary linearly-damped equation of motion --- genuine viscous friction,
    obtained from the stationarity of a single functional, with no dissipation
    function added by hand. This is exactly the kind of open-system equation
    discussed elsewhere in this book (Meshchersky, Bateman's dual Lagrangian,
    Abraham--Lorentz); Herglotz's principle is the general variational scaffold
    those examples are individual instances of.
\end{example}

Herglotz gave this construction in unpublished 1930 Göttingen lectures; it
re-entered the literature only via Georgieva and Guenther's rediscovery (2002),
and a decade later Bravetti, Cruz and Tapias recognized it as the variational
principle native to \emph{contact}, rather than symplectic, geometry --- the
natural geometric home for dissipative mechanics, exactly as symplectic
geometry is the natural home for the conservative principles of
Sections~\ref{sec:hamilton}--\ref{sec:hertz}. This is a second illustration of
Section~\ref{sec:primacy}'s thesis, running in the opposite direction:
friction was never in doubt as physics, and it took
decades to notice that a single genuine variational principle --- not a
patched-on dissipation function --- could be engineered to produce it.

\section{Schwinger's action principles}
\label{sec:schwinger}

Every principle so far, once phase space is introduced, treats $p_i$ as
\emph{derived}: $p_i:=\partial L/\partial\dot q^i$, fixed in terms of $q,\dot
q$ before any variation is taken. There is an older and logically different
construction, due to Cartan (1922) and extended to field theory by Weiss (1936,
1938), in which $q^i(t)$ and $p_i(t)$ are varied as \emph{independent}
functions, with no relation between them assumed in advance. Schwinger opens
his 1951 paper on this classical construction before promoting it to operators.

\begin{definition}
\textbf{The classical phase-space action principle (Cartan--Weiss--Schwinger).}
    Let $q^i(t)$ and $p_i(t)$ be independent trial functions on $[t_1,t_2]$,
    and define
\[
\tilde S[q(\cdot),p(\cdot)] := \int_{t_1}^{t_2}\big(p_i\dot q^i - H(q,p,t)\big)\,dt .
\]
The true motion renders $\tilde S$ stationary under independent variations
    $\delta q^i$ (vanishing at $t_1,t_2$) and $\delta p_i$ (unconstrained at
    the endpoints, since $p$ appears undifferentiated).
\end{definition}

Independent stationarity in $\delta p_i$ and $\delta q^i$ yields, separately,
\[
\dot q^i = \frac{\partial H}{\partial p_i}
\qquad\text{and}\qquad
\dot p_i = -\frac{\partial H}{\partial q^i},
\]
i.e.\ \emph{both} of Hamilton's equations, with no leftover boundary term
beyond $\big[p_i\,\delta q^i\big]_{t_1}^{t_2}$. Indeed,
\[
\delta\tilde S = \int_{t_1}^{t_2}\left[\left(\dot q^i - 
\frac{\partial H}{\partial p_i}\right)\delta p_i - \left(\dot p_i + 
\frac{\partial H}{\partial q^i}\right)\delta q^i\right]dt + 
\big[p_i\,\delta q^i\big]_{t_1}^{t_2},
\]
obtained by integrating the $p_i\delta\dot q^i$ term by parts. Since $\delta
p_i$ and $\delta q^i$ are independent and unconstrained in the interior, the
two bracketed coefficients must vanish separately, and $\delta q^i(t_1)=\delta
q^i(t_2)=0$ kills the boundary term.

This is a different logical achievement from Hamilton's principle of
Section~\ref{sec:hamilton}. There, only $q$ is varied and $p:=\partial
L/\partial\dot q$ is a definition, so only the second Hamilton equation is
derived --- the first is built into the definition of $p$ before variation ever
starts. Here neither equation is assumed: treating $q$ and $p$ as coordinates
of equal standing on phase space, with no a priori link between $p$ and $\dot
q$, derives \emph{both} canonical equations symmetrically. It is exactly this
symmetry between $q$ and $p$ --- and the resulting boundary term $p_i\,\delta
q^i$, free of any assumed relation between the two --- that Schwinger promotes
to an operator statement next.

\section{Marinov's phase action}
\label{sec:marinov-phase-action}

Hamilton--Jacobi theory demands more than solving a partial differential
equation: constructing Hamilton's principal function $\mathcal A(Q,q_0;t)$ from
\[
\partial\mathcal A/\partial t + H(Q,\partial\mathcal A/\partial Q;t)=0
\]
requires a \emph{complete integral} --- an $f$-parameter family of solutions
--- not merely an initial-value problem. And because the two-point
boundary-value problem $q(0)=q_0$, $q(t)=Q$ generically has several solutions
(Section~\ref{sec:nonequivalence}), $\mathcal A$ is generically
\emph{multivalued}. Marinov (1979) replaced $\mathcal A$ with a genuinely
different object, defined on phase space rather than on pairs of
configurations, for which neither defect occurs.

\begin{definition}
\textbf{Marinov's phase action.} For a trajectory $x(\tau)=(q(\tau),p(\tau))$
    solving Hamilton's equations $\dot x=\omega\nabla H$ with $x(0)=\xi$, let
    $z:=\tfrac12(x(t)+\xi)$ be the phase-space \emph{midpoint} of its
    endpoints. The \emph{phase action} $\Phi(z;t)$ is defined implicitly by
\[
x(t) - \xi = \omega\,\nabla_z\Phi(z;t) .
\]
\end{definition}

$\Phi$ is the unique solution of the Cauchy problem
\[
\frac{\partial\Phi}{\partial t} = H\!\left(z+\tfrac12\nabla\Phi;\,t\right),
\qquad \Phi(z;0)=0 .
\]
Unlike the Hamilton--Jacobi equation for $\mathcal A$, this is a genuine
\emph{initial-value} problem: $\Phi(z;0)=0$ pins down a unique solution by the
Cauchy--Kovalevskaya theorem, with no complete integral to hunt for. $\Phi$ is
therefore single-valued on phase space, in sharp contrast to $\mathcal A$.

Writing $z=(r,k)$ for the configuration and momentum parts of the midpoint,
\[
\mathcal A\!\left(r+\tfrac12\rho,\ r-\tfrac12\rho;\,t\right) = k\rho -
\Phi(r,k;t), \qquad \rho=\partial\Phi/\partial k .
\]

This identifies exactly where the multiplicity of
Section~\ref{sec:nonequivalence} comes from: a Legendre transform of a
perfectly well-behaved, single-valued function can be, and generically is,
multivalued --- the same mechanism that turns one thermodynamic potential into
several branches of another. The bifurcations in $\mathcal A$ are an artifact
of which variable ($\rho$, tied to the two separate endpoints) is used to label
solutions, not a pathology in the underlying dynamics, which $\Phi$ describes
without ambiguity.

On the true trajectory,
\[
\Phi(z;t) = \int_0^t\Big[H(X) - \tfrac12(X\circ\dot X)\Big]\,d\tau + (\xi\circ z), \qquad \xi\circ z:=\tilde\omega^{kl}\xi_kz_l,
\]
where $X(\tau)$ ranges over trajectories subject only to the \emph{midpoint}
boundary condition $X(0)+X(t)=2z$ --- fixing not either endpoint individually,
but their phase-space average.

This is a fourth kind of boundary condition for a classical variational
principle, alongside Hamilton's (fix both configuration endpoints),
Maupertuis'/GMP's (fix the energy), and Gauss's/Hertz's (fix nothing beyond the
current instant). As with Hamilton's principle, the second variation of this
functional has no definite sign: the true trajectory is a saddle point, never a
minimum or maximum, among curves sharing its midpoint.

The isotropic oscillator's phase action blows up at $\kappa t=\pi$ and every
subsequent half-period, marking a genuinely new kind of singularity, distinct
from the conjugate points of Section~\ref{sec:jacobi}: at $\kappa t=\pi$ the
oscillator sends every $\xi$ to $-\xi$, so the midpoint
$z=\tfrac12(x(t)+\xi)=0$ regardless of $\xi$ --- every trajectory shares the
same midpoint, the map trajectory$\mapsto z$ stops being invertible, and $\Phi$
cannot remain a well-defined function of $z$ alone.

\section{Summary and classification}
\label{sec:summary}

Table~\ref{tab:summary} cross-classifies every principle of this chapter along
the same three axes used throughout: what is varied, what is held fixed versus
left free, and whether the comparison is a differential (single-instant or
single-operator) or an integral (whole-history) one.

\begin{table}[ht]
\centering
\footnotesize
\renewcommand{\arraystretch}{1.35}
\begin{tabularx}{\textwidth}{@{}lXXXl@{}}
\toprule
Principle & Object varied & Held fixed & Free & Cl./qu. \\
\midrule
Fermat & $\int n\,ds$ & endpoints & path shape & optics \\
Hamilton (HP) & $S=\int L\,dt$ & $q_1,q_2,t_1,t_2$ & --- & classical \\
Hamilton, reciprocal (RHP) & $T$ & $q_1,q_2,S$ & $t_1,t_2$ & classical \\
Maupertuis, textbook (MP) & $W=\int p\,dq$ & $q_1,q_2,E$ pointwise & transit time & classical \\
Maupertuis, generalized (GMP) & $W=\int p\,dq$ & $q_1,q_2,\bar E$ & transit time & classical \\
Maupertuis, reciprocal (RMP) & $\bar E$ & $q_1,q_2,W$ & transit time & classical \\
Jacobi & $\int ds_J$ & $q_1,q_2,E$ & transit time & geometrized \\
Gauss & $Z(\ddot q)$ & $q,\dot q$, instant & accelerations & differential \\
Hertz & curvature & $q,\dot q$, instant, $V\!\equiv\!0$ & accelerations & differential \\
Herglotz & $z(t_2)$ & $q_1,q_2,z(t_1)$ & --- & classical, dissipative \\
Schwinger, classical & $\tilde S[q,p]$ & $q_1,q_2$ & all of $p(t)$ & phase-space \\
Marinov, phase action & $\int[H-\tfrac12(X\circ\dot X)]d\tau$ & midpoint $z=\tfrac12(X(0){+}X(t))$, $t$ & --- & classical \\
\bottomrule
\end{tabularx}
\caption{A cross-classification. ``Differential'' means the principle compares
    alternatives at a single instant or as a single operator relation;
    ``integral'' (the rest) means it compares whole trial histories over a
    finite stretch.}
\label{tab:summary}
\end{table}


\section{Extension to fields}

Everything above generalizes, with $q^i(t)$ replaced by a field $\varphi^a(x)$
on spacetime, by promoting Hamilton's principle to
\[
\delta\!\int L(\varphi,\partial_\mu\varphi)\,d^4x=0,
\]
giving the field Euler--Lagrange equations
\[
\partial_\mu\!\left(\frac{\partial L}{\partial(\partial_\mu\varphi^a)}\right)
= \frac{\partial L}{\partial\varphi^a}.
\]
The Hamilton/Maupertuis distinction --- which coordinate plays the role of
``time'' to be fixed versus freed --- becomes, in field theory, the question of
treating one spacetime direction asymmetrically (ordinary Hamiltonian field
theory) versus all of them symmetrically (the De Donder--Weyl covariant
Hamiltonian formalism and multiple-integral calculus of variations, treated as
their own chapter of this book). Gauss's and Hertz's differential principles do
not obviously survive the passage to fields in a useful form, which is itself
worth noting: instantaneous, finite-dimensional least-constraint reasoning does
not have an established infinite-dimensional counterpart of comparable power.

\section{Equations of motion are primary}
\label{sec:primacy}

None of the principles in this chapter is more ``true'' than the others in the
sense of better describing the real-world trajectory: for a holonomic,
energy-conserving system they are jointly compatible, and physically
indistinguishable in their predictions. What differs is which real-world
question each idealization is built to answer directly. ``What configuration is
the system in at time $t$?'' is answered most directly by Hamilton's principle.
``What is the geometric shape of the orbit, independent of how fast it is
traversed?'' is answered most directly by Maupertuis' and Jacobi's. ``What
happens to constraint forces at this instant?'' is Gauss's question. ``What are
the possible outcomes of a measurement, and with what amplitude?'' is
Schwinger's and Feynman's. Section~\ref{sec:nonequivalence} was the cautionary
half of that remark: the compatibility just described holds trajectory by
trajectory, not problem by problem.

Put the two observations together and a sharper claim follows, and this chapter
is the argument for it.

\begin{motto}
\textbf{Thesis.} The equations of motion are the real-world content of a
    mechanical theory: the differential, causal, initial-value statement that
    is actually checked against experiment and that alone stays invariant
    across every reformulation in this chapter. A variational principle is an
    ideal-world heuristic for \emph{finding} those equations --- a scaffold
    engineered so that stationarity reproduces a differential law already fixed
    by other means. It is not itself a law of nature, and no metaphysical
    content should be read into the word ``least.''
\end{motto}

\begin{remark}[underdetermination]
Table~\ref{tab:summary} lists objects varied of genuinely different
    mathematical kinds --- a time integral ($S$), a curve integral ($W$), a
    pointwise quadratic form ($Z(\ddot q)$), an operator matrix element --- all
    of which are stationary exactly on the solutions of the same equations of
    motion. If any one of these were the law nature ``obeys,'' the others would
    need explaining away; there is no principled way to prefer $\int L\,dt$
    over $Z(\ddot q)$ as the fundamental fact and demote the other to derived
    curiosity. What is invariant across all of them, and hence what deserves to
    be called primary, is the one thing they all agree on: the equations of
    motion themselves.
\end{remark}

\begin{remark}[Gauss and Hertz remove the teleology]
A global, fixed-endpoint statement like Hamilton's principle can sound as
    though the system ``surveys'' every possible history between two moments
    before choosing the actual one --- historically the most seductive (and
    most criticized) reading of least action. Gauss's principle removes this
    appearance entirely:  it is 
    equivalent, instant by instant, to $(m_a\ddot{\vec r}_a-\vec
    F_a)\cdot\delta\vec r_a=0$, i.e.\ to Newton's second law together with
    d'Alembert's principle, with no reference to the future or to any global
    comparison of histories. The same content that Hamilton's principle
    packages as ``compare the whole path'' is available, undiminished, as a
    purely local, instant-by-instant fact. Hertz's principle
    (Section~\ref{sec:hertz}) is the same point in geometric dress: force-free
    motion is just the straightest available path, an entirely local notion of
    curvature. Nothing about the physics requires the global packaging; the
    global packaging is a choice of heuristic, not a discovery about how nature
    decides.
\end{remark}

\begin{remark}[fragility shows the direction of justification]
Naively extremizing $\int L\,dt$ under a nonholonomic
    constraint gives the \emph{wrong} equations of motion (the vakonomic ones);
    the \emph{right} equations (Lagrange--d'Alembert) are fixed independently,
    by direct mechanical reasoning about constraint forces and virtual work,
    and it is against that independently-known target that the naive
    variational guess is judged and rejected. Section~\ref{sec:gauss} was
    accepted as ``the right principle'' \emph{because} it reproduces the
    equations of motion known to be correct on other grounds --- not the other
    way around. If a variational principle were itself the primary law, it
    would never need to be repaired against a differential equation assumed to
    be more trustworthy than it is.
\end{remark}

\begin{remark}[the equations of motion are the robust object]
Section~\ref{sec:nonequivalence} showed that once a potential has several
    minima, the variational \emph{problem} can fragment: the textbook
    Maupertuis competition can have an empty domain, the Generalized Maupertuis
    Principle's competition can fail to contain an extremal, and brake-orbit
    multiplicity can make ``the'' fixed-energy solution ambiguous. Through all
    of this the equation of motion itself --- an ordinary differential equation
    with a well-posed initial-value problem --- never fragments, never loses a
    solution, and never multiplies one; it is exactly as well behaved for ten
    wells as for one. The variational reformulation is the fragile,
    construction-dependent object; the differential law is the robust one.
\end{remark}

\begin{remark}[even ``derived versus assumed'' is bookkeeping]
Section~\ref{sec:schwinger} made a further, sharper version of the
    same point at the classical level: whether $\dot q^i=\partial H/\partial
    p_i$ counts as a consequence of a variational principle or as a definition
    adopted before any variation is taken depends entirely on whether one
    varies $q$ alone (ordinary Hamilton's principle) or $q$ and $p$
    independently (the phase-space principle). The two canonical equations
    themselves do not change; only the bookkeeping of which one gets called
    ``derived'' does. A fact about nature should not depend on which
    variational scaffold was used to reach it.
\end{remark}

\begin{remark}[the name itself overclaims]
Even calling these ``least'' action is often simply false. Once a trajectory
    passes a conjugate point --- exactly where the curvature of the Jacobi
    metric of Section~\ref{sec:jacobi} turns the second variation indefinite
    --- the stationary trajectory is a \emph{saddle} of $S$, not a minimum.
    Hamilton himself never claimed a minimum principle, only a stationary one;
    ``least action'' is a historical label surviving from a reading the
    mathematics does not support in general.
\end{remark}

The historical record backs this reading of the founders' own intentions, not
merely a modern gloss on them. Maupertuis' 1746 memoir was titled, without
irony, ``the laws of motion and of rest deduced from a metaphysical principle''
--- he offered least action as evidence of a providential economy in nature,
not as a mathematical convenience. But d'Alembert, Lagrange, Hertz and Jacobi
--- three of whom are principal subjects of this very chapter --- explicitly
rejected any ontological economy-of-nature behind their own principles,
treating least action as what Panza and Goldstine call a figurative scientific
model rather than a metaphysical discovery. Mach's \emph{Science of Mechanics}
(1883) makes the case most sharply, in a chapter pointedly titled
\emph{Theological, Animistic, and Mystical Points of View in Mechanics},
followed immediately by one titled \emph{The Economy of Science}: for Mach,
least-action-type statements are \emph{theorems} --- consequences, economically
redescribing what is already contained in Newton's laws --- not
\emph{principles} in the sense of directly-intuited primitive facts (his
examples of genuine principles are the lever and virtual displacements). Their
teleological gloss, Mach argued, was a product of a theologically-tempered era,
not a discovery about how nature operates. Helmholtz, writing in the same
period, offered the more constructive half of the same demotion: whatever their
metaphysical status, least-action-type statements had proved, again and again,
a powerful heuristic and leitmotif for finding the governing equations of new
domains of physics before those equations were otherwise known --- exactly the
role this book assigns them, and exactly why this chapter surveys so many of
them side by side rather than picking one as ``the'' principle.


